\documentclass[10pt,a4paper]{amsart}
\usepackage[margin=19mm]{geometry}
\usepackage{amsmath,amssymb,mathtools}
\usepackage[scr=rsfs,cal=euler]{mathalfa}
\usepackage{xfrac}
\usepackage[unicode,hidelinks,bookmarks=false]{hyperref}
\newcommand{\ie}{i.e.\ }
\newcommand{\cf}{cf.\ }
\newcommand{\resp}{respectively\ }
\newcommand{\Subsection}{\subsection}
\providecommand{\nobreakdash}{\nobreak\hbox{-}}
\setlength{\emergencystretch}{2em}
\multlinegap0pt
\usepackage{bm}
\usepackage{bbm}
\usepackage{dsfont}
\usepackage{xspace}
\usepackage{cancel}
\usepackage{mathtools}
\usepackage{amsthm}
\makeatletter
\@ifundefined{theo}{\newtheorem{theo}{Theo}}{}
\@ifundefined{gathered}{\newtheorem{gathered}{Gathered}}{}
\@ifundefined{lemm}{\newtheorem{lemm}[theo]{Lemma}}{}
\makeatother
\newcommand{\CD}{\mathcal{D}}
\newcommand{\CU}{\mathcal{U}}
\newcommand{\C}{\mathbb{C}}
\newcommand{\R}{\mathbb{R}}
\newcommand{\T}{\mathbb{T}}
\newcommand{\dif}{\mathrm{d}}
\newcommand{\eu}{\mathrm{e}}
\newcommand{\inv}{^{-1}}
\newcommand{\wh}{\widehat}
\newcommand{\wt}{\widetilde}
\renewcommand{\Re}{\operatorname{Re}}
\renewcommand{\epsilon}{\varepsilon}
\title[Spectral rigidity of two-dimensional Liouville tori]{Spectral rigidity of two-dimensional Liouville tori}
\author{Joscha Henheik, Vadim Kaloshin, Yunzhe Li, Amir Vig}
\date{Source: arXiv:2511.10398v2 (CC BY 4.0)}
\begin{document}
\maketitle

\noindent\emph{Source and licence.} Spectral rigidity of two-dimensional Liouville tori. Version arXiv:2511.10398v2 (CC BY 4.0), distributed under the Creative Commons Attribution 4.0 International licence, as recorded in the arXiv licence metadata for this exact version and shown on the abstract page https://arxiv.org/abs/2511.10398v2. Full rights notice: licenses/arxiv-2511.10398-CC-BY-4.0.txt.\par\smallskip

\noindent\textit{Editorial scope.} Counted unit: the Lemm whose source label is \texttt{lem:Gam determines f12} in the article (source0.tex lines 2584--2587, source label \texttt{lem:Gam determines f12}), delivered together with the article's own complete original proof of it (source0.tex lines 2588--2691, one \texttt{proof} environment of 104 lines). No mathematics is written, repaired or re-derived: statement and proof are the article's own text line for line. Required context retained uncounted: eq:Gam ends (lines 2560--2565); eq:layer-cake J (lines 2577--2579). Every definition, display and label of those lines is retained unchanged. Omitted from this package and not redistributed: 1--2559, 2566--2576, 2580--2583, 2692--3099. Nothing inside a retained excerpt is omitted or altered: every retained body is delivered exactly as the article's lines are written, so no omission receipt of any kind accompanies this package. Citations: the retained excerpts carry no citation command at all, so no bibliography entry of the article is transcribed and no reference is redistributed. Reference adaptation: none was needed and none was made; every reference inside the delivered unit resolves inside it, so no unresolved reference or citation is produced. Printed slips kept literal: no printed word, symbol, hypothesis or number of the counted statement or of its proof was altered. Layout and class: the article's own class is \texttt{amsart} with packages including amsmath, amsthm, amssymb, graphics, color, hyperref, xcolor, babel, tikz, tikz-cd, geometry, mathrsfs, soul, subcaption; the wrapper is the lane's pinned \texttt{amsart} editorial wrapper at 10pt on A4, which restates the theorem-like environments the retained text needs and adds the article's own macro definitions that the retained text needs (\texttt{\textbackslash C}, \texttt{\textbackslash CD}, \texttt{\textbackslash CU}, \texttt{\textbackslash R}, \texttt{\textbackslash Re}, \texttt{\textbackslash T}, \texttt{\textbackslash dif}, \texttt{\textbackslash epsilon}, \texttt{\textbackslash eu}, \texttt{\textbackslash inv}, \texttt{\textbackslash wh}, \texttt{\textbackslash wt}), with extra wrapper packages (bm, bbm, dsfont, xspace, cancel, mathtools). The delivered numbering is the unit's own. Assets: no retained span contains a figure environment, a graphics-inclusion command, a PDF-page inclusion, a table or a TikZ picture; no image file is copied into this package, the package declares no asset dependency and no image file is redistributed with it. Licence: version arXiv:2511.10398v2 of this article is distributed under the Creative Commons Attribution 4.0 International licence, as recorded in the arXiv licence metadata for this exact version and shown on the abstract page https://arxiv.org/abs/2511.10398v2. A full rights notice is delivered at licenses/arxiv-2511.10398-CC-BY-4.0.txt. Characters: every retained excerpt is pure ASCII, so the delivered engine, XeLaTeX with its default Latin Modern text font, reports no missing character.
\begin{equation}\label{eq:Gam ends}
\begin{gathered}
    \Big(\int_\T \sqrt{f_1(x) - \min f_1} \dif x,\ \int_\T \sqrt{1 + \min f_1 + f_2(x)} \dif x\Big) \quad \text{and}\\
    \Big(\int_\T \sqrt{1 + \min f_2 + f_1(x)} \dif x,\ \int_\T \sqrt{f_2(x) - \min f_2} \dif x \Big).
\end{gathered}
\end{equation}

\begin{equation}\label{eq:layer-cake J}
	J_f(e) = \frac{1}{2} \int_{-e}^\infty \frac{L_f(t)}{\sqrt{t + e}} \dif t = \int_{-e}^\infty \sqrt{t + e} \;\dif \mu_f(t).
\end{equation}

\begin{lemm}\label{lem:Gam determines f12}
	Let $(\wt f_1, \wt f_2)$ and $(f_1, f_2)$ be two pairs of $C^2$ functions $\T \to \R$ with $\min f_1 + \min f_2 > -1$ and $\min \wt f_1 + \min \wt f_2 > -1$.
	Then $\Gamma_{f_1, f_2} = \Gamma_{\wt f_1, \wt f_2}$ if, and only if there exists $c \in \R$ such that $\wt f_1$ is a rearrangement of $f_1 + c$ and $\wt f_2$ is a rearrangement of $f_2 - c$.
\end{lemm}

\begin{proof}
	We have observed the ``if'' part. 
	For the ``only if'' part, suppose $\Gamma_{\wt f_1, \wt f_2} = \Gamma_{f_1, f_2}$ as subsets of $\R^2$ for some $C^2$ functions $f_1, f_2, \wt f_1$ and $\wt f_2$ satisfying the conditions of the lemma.
	Modulo transformations of the type $(f_1, f_2) \mapsto (f_1 + c, f_2 - c)$, we may assume that $\min f_1 = \min \wt f_1 = 0$.
	To show the lemma it is enough to prove that $\wt f_j$ is a rearrangement of $f_j$ for $j = 1,2$.

\medskip

    We begin by matching the endpoints \eqref{eq:Gam ends} of the curves $\Gamma_{\wt f_1, \wt f_2}$ and $\Gamma_{f_1, f_2}$, which in particular implies
    \begin{equation*}        \int_{\T} \sqrt{f_j(x) - \min f_j}\ \dif x = 
        \int_{\T} \sqrt{\wt f_j(x) - \min \wt f_j}\ \dif x \qquad (j = 1,2).
    \end{equation*}
	Since $J_f(e)$ is real-analytic and strictly increasing for $e > -\min f$, it follows that 
    \begin{equation*}        J_{\wt f_j}\inv \circ J_{f_j} : (-\min f_j, \infty) \to  (-\min \wt f_j, \infty) \qquad (j = 1,2)
    \end{equation*}
    are real-analytic diffeomorphisms.
    Define 
    \begin{equation*}    \varphi(e) :=  \begin{cases}
            J_{\wt f_1}\inv \circ J_{f_1} (e) & (e > -\min f_1) \\
            1 - J_{\wt f_2}\inv \circ J_{f_2} (1 - e)  & (e < 1 + \min f_2).
        \end{cases}
    \end{equation*}
    Since $\Gamma_{f_1, f_2}$ projects injectively to $x_j$-coordinate for $j = 1,2$, we can verify that the two definitions of $\varphi(e)$ are consistent for $-\min f_1 < e < 1 + \min f_2$ and
    \begin{equation}\label{eq:J reparam}
        (J_{f_1}(e), J_{f_2}(1- e)) = (J_{\wt f_1}(\varphi(e)), J_{\wt f_2}(1- \varphi(e))) \qquad  (-\min f_1 < e < 1 + \min f_2).
    \end{equation}
    By analyticity of $J_{f_j}$ and $J_{\wt f_j}$ over their respective domains, we deduce that $\varphi: \R \to \R$ is a real-analytic diffeomorphism.
    We remark that $\varphi(1 + \min f_2) = 1 + \min \wt f_2$ and 
    \begin{equation*}    
        \varphi(0) = \varphi(-\min f_1) = - \min \wt f_1 = 0
    \end{equation*} 
 since we have assumed $\min f_1 = \min \wt f_1 = 0$.

\medskip

	Let us denote $f = f_1$ and $\wt f = \wt f_1$. Consider
	\begin{equation*}
		\wh J_{f}(e) = \int_0^\infty \sqrt{t + e} \; \dif \mu_f(t).
	\end{equation*}
	Choosing the principal branch of the square root, we define $\wh J_f$ analytically on the domain $\C \setminus (-\infty, 0]$.
	Observe that $J_f|_{(0, \infty)}$ is exactly the restriction of $\wh J_f$ on $(0, \infty)$.
	Hence $J_f$ determines $\wh J_f$  by analytic continuation.

\medskip

	Let $\mathbb{H}^\pm$ denote the upper/lower complex half-plane. 
	If $e \in \mathbb{H}^\pm$, then $\sqrt{t + e}$ lies in the first/fourth quadrant of $\C$ for any $t > 0$. It follows from convexity that $\wh J_f$ maps $\mathbb{H}^\pm$ into the first/fourth quadrant respectively.
	In fact, we show that $\wh J_f$ restricts to biholomorphisms
	\begin{equation}\label{eq:hat J bihol}
		\wh J_f: \mathbb{H}^\pm \xrightarrow{\sim} \CD^\pm,
	\end{equation}
	where $\CD^\pm$ are the unbounded domains in the first/fourth quadrant of $\C$ bounded by the curve 
	\begin{equation}\label{eq:D boundary}
		e \mapsto \lim_{\epsilon \to 0^\pm}\wh J_f (e + i \epsilon) = \int_{-e}^{\max f} \sqrt{t + e} \; \dif \mu_f(t) \pm i \int_{0}^{-e} \sqrt{|t + e|}\; \dif \mu_f(t).
	\end{equation}
	Notice that \eqref{eq:D boundary} parametrizes the image of the boundary of $\mathbb{H}^\pm$ under $\wh J_f$, and it is understood that the first integral is zero if $-e > \max f$ and similarly the second integral is zero for $-e < 0$.

\medskip

	To show \eqref{eq:hat J bihol}, we first observe that, since $\mu_f$ is a finite measure supported in the interval $[0, \max f]$, we have 
	\begin{equation}\label{eq:hat J asymptotic}
		|z|^{-1/2} |\wh J_f(z)|\to 1 \quad \text{uniformly as $|z| \to \infty$}
	\end{equation}
	Let $z_0 \in \CD^+$.
	Then, for $R$ sufficiently large, the point $z_0$ is enclosed by a closed loop formed by the image of the semicircle $\{z \in \mathbb{H}^+ : |z| = R\}$ under the function $\wh J_f$ and the boundary \eqref{eq:D boundary}.
	By shrinking $R$, the usual argument using contractible loops shows that $z_0 \in \wh J_f(\mathbb{H}^+)$. Hence $\wh J_f(\mathbb{H}^+) = \CD^+$.
	Similarly $\wh J_f(\mathbb{H}^-) = \CD^-$.
	To see the injectivity of $\wh J_f$, suppose $\wh J_f(z_1) = \wh J_f(z_2)$. Then
	\begin{equation*}
		0 = \wh J_f(z_1) -  \wh J_f(z_2) = \int_0^\infty \sqrt{t + z_1} - \sqrt{t + z_2} \dif \mu_f(t) = (z_1 - z_2) \int_0^\infty \frac{\dif \mu_f(t)}{\sqrt{t + z_1} + \sqrt{t + z_2}}.
	\end{equation*}
	Since $\Re \sqrt{t + z} > 0$ for all $t > 0$ and $z \in \C \setminus (-\infty, 0]$, it is easy to see that the last integral above does not vanish. It follows that $z_1 = z_2$ and hence $\wh J_f$ is injective.
	It is also clear that $\wh J_f$ has nonvanishing derivative everywhere. It follows that \eqref{eq:hat J bihol} are indeed biholomorphisms.
	Similarly, we have $\wh J_{\wt f}: \mathbb{H}^\pm \xrightarrow{\sim} \wt{\CD}^\pm$ where $\wt{\CD}^\pm$ are domains defined analogously with boundaries given by \eqref{eq:D boundary} with $f$ replaced by $\wt f$.
	
\medskip

	Recall that we have $J_f(e) = J_{\wt f}(\varphi(e))$ for $e > 0$ and $\varphi: \R \to \R$ real-analytic.
	Extending $\varphi$ analytically over a complex neighborhood $\CU \supset \R$, we have $\wh J_f = \wh J_{\wt f} \circ \varphi$ over $\CU \setminus (-\infty, 0]$.
	Shrinking $\CU$ if necessary we may assume $\varphi(\CU \cap \mathbb{H}^\pm) \subset \mathbb{H}^\pm$.
	It follows by taking the limit $\varphi(e + i \epsilon)\to \varphi(e)$ for $e \in \R$ that 
	\begin{equation*}		\lim_{\epsilon \to 0^\pm} \wh J_f(e + i \epsilon) = \lim_{\epsilon \to 0^\pm} \wh J_{\wt f}(\varphi(e) + i \epsilon).
	\end{equation*}
	Consequently, the boundary curves \eqref{eq:D boundary} defined by $f$ and by $\wt f$ are identical up to a reparametrization by $\varphi$ and hence both $\wh J_{\wt f}$ and $\wh J_f$ are biholomorphisms $\mathbb{H}^\pm \xrightarrow{\sim} \CD^\pm$ onto the same image.
	We can therefore extend the original definition of $\varphi$ to an analytic function $\varphi = \wh J_{\wt f}\inv \circ \wh J_f$  over $\C\setminus (-\infty, 0]$.
	Since $\varphi$ is real-analytic over $\R$, by the Schwarz reflection principle, it is an entire function.

\medskip

	Since $\varphi(0) = 0$, the function $z \mapsto \varphi(z) /z$ is entire and bounded by the asymptotics \eqref{eq:hat J asymptotic}.
	It then follows from Liouville's theorem that $\varphi(z) = az$ for some $a \in \R$.
	But \eqref{eq:hat J asymptotic} forces $a = 1$. Thus $\varphi(z) = z$ and we obtain accordingly $J_f = J_{\wt f}$ and $J_{f_2}(1 - e) = J_{\wt f_2}(1 - e)$. Since $L_f(e) = L_{\wt f}(e) = 1$ for $e < 0$, Equation \eqref{eq:layer-cake J} implies 
	\begin{equation*}		\int_{0}^\infty \frac{L_f(t)}{\sqrt{t + e}} \dif t = \int_{0}^\infty \frac{L_{\wt f}(t)}{\sqrt{t + e}} \dif t \quad \text{for all $e > 0$}.
	\end{equation*}
	Using the identity $\int_0^\infty s^{-1/2}\eu^{-(t + e) s} \dif s = \pi^{1/2} (e + t)^{-1/2}$, we have 
	\begin{equation*}
		\int_{0}^\infty \frac{L_f(t)}{\sqrt{t + e}} \dif t = \frac{1}{\sqrt{\pi}} \int_{0}^\infty \int_0^\infty  L_f(t) \frac{\eu^{-(e + t)s}}{\sqrt{s}} \dif s \dif t = \frac{1}{\sqrt{\pi}} \mathscr{L}_s \bigg[  \frac{\mathscr{L}_t[L_f(t)](s)}{\sqrt{s}}\bigg](e)
	\end{equation*}
	where $\mathscr{L}_t$  and  $\mathscr{L}_s$  denote the Laplace transform in the variable $t$ and $s$ respectively.
	By injectivity of Laplace transform, we have $L_f = L_{\wt f}$.
	By the same argument applied to $J_{f_2}(1 - e) = J_{\wt f_2}(1 - e)$ for $e < 1 + \min f_2$, we deduce $L_{f_2} = L_{\wt f_2}$. 
	Hence $\wt f_j$ is a rearrangement of $f_j$ as desired. 
	The proof is complete.
\end{proof}

\end{document}
